Networking met Linux is soms lastig. We kunnen met Linux een eenvoudige ethernet-kaart aansturen voor de verbinding met een LAN, maar er kan ook ADSL mee aangestuurd worden als Linux ingezet wordt op een router naar het Internet. De meeste apparaten die je thuis-netwerk aan het Internet koppelen zijn bevatten een Linux systeem. Linux kan ook gebruikt worden om een router of een switch te bouwen. Al deze functionaliteit kan al gauw onoverzichtelijk worden. Gelukkig zijn er veel tools die het leven makkelijker maken. In dit document zullen we de basis van Linux netwerken beschrijven. We gaan ervan uit dat je een desktop, laptop of server hebt die je aan het netwerk wilt koppelen. We zullen je dus laten zien hoe je ethernet, wifi en IP configureert.

